% Options for packages loaded elsewhere
\PassOptionsToPackage{unicode}{hyperref}
\PassOptionsToPackage{hyphens}{url}
%
\documentclass[
  12pt,
]{article}
\usepackage{amsmath,amssymb}
\usepackage{iftex}
\ifPDFTeX
  \usepackage[T1]{fontenc}
  \usepackage[utf8]{inputenc}
  \usepackage{textcomp} % provide euro and other symbols
\else % if luatex or xetex
  \usepackage{unicode-math} % this also loads fontspec
  \defaultfontfeatures{Scale=MatchLowercase}
  \defaultfontfeatures[\rmfamily]{Ligatures=TeX,Scale=1}
\fi
\usepackage{lmodern}
\ifPDFTeX\else
  % xetex/luatex font selection
\fi
% Use upquote if available, for straight quotes in verbatim environments
\IfFileExists{upquote.sty}{\usepackage{upquote}}{}
\IfFileExists{microtype.sty}{% use microtype if available
  \usepackage[]{microtype}
  \UseMicrotypeSet[protrusion]{basicmath} % disable protrusion for tt fonts
}{}
\makeatletter
\@ifundefined{KOMAClassName}{% if non-KOMA class
  \IfFileExists{parskip.sty}{%
    \usepackage{parskip}
  }{% else
    \setlength{\parindent}{0pt}
    \setlength{\parskip}{6pt plus 2pt minus 1pt}}
}{% if KOMA class
  \KOMAoptions{parskip=half}}
\makeatother
\usepackage{xcolor}
\usepackage[margin=2.5cm]{geometry}
\usepackage{color}
\usepackage{fancyvrb}
\newcommand{\VerbBar}{|}
\newcommand{\VERB}{\Verb[commandchars=\\\{\}]}
\DefineVerbatimEnvironment{Highlighting}{Verbatim}{commandchars=\\\{\}}
% Add ',fontsize=\small' for more characters per line
\usepackage{framed}
\definecolor{shadecolor}{RGB}{248,248,248}
\newenvironment{Shaded}{\begin{snugshade}}{\end{snugshade}}
\newcommand{\AlertTok}[1]{\textcolor[rgb]{0.94,0.16,0.16}{#1}}
\newcommand{\AnnotationTok}[1]{\textcolor[rgb]{0.56,0.35,0.01}{\textbf{\textit{#1}}}}
\newcommand{\AttributeTok}[1]{\textcolor[rgb]{0.13,0.29,0.53}{#1}}
\newcommand{\BaseNTok}[1]{\textcolor[rgb]{0.00,0.00,0.81}{#1}}
\newcommand{\BuiltInTok}[1]{#1}
\newcommand{\CharTok}[1]{\textcolor[rgb]{0.31,0.60,0.02}{#1}}
\newcommand{\CommentTok}[1]{\textcolor[rgb]{0.56,0.35,0.01}{\textit{#1}}}
\newcommand{\CommentVarTok}[1]{\textcolor[rgb]{0.56,0.35,0.01}{\textbf{\textit{#1}}}}
\newcommand{\ConstantTok}[1]{\textcolor[rgb]{0.56,0.35,0.01}{#1}}
\newcommand{\ControlFlowTok}[1]{\textcolor[rgb]{0.13,0.29,0.53}{\textbf{#1}}}
\newcommand{\DataTypeTok}[1]{\textcolor[rgb]{0.13,0.29,0.53}{#1}}
\newcommand{\DecValTok}[1]{\textcolor[rgb]{0.00,0.00,0.81}{#1}}
\newcommand{\DocumentationTok}[1]{\textcolor[rgb]{0.56,0.35,0.01}{\textbf{\textit{#1}}}}
\newcommand{\ErrorTok}[1]{\textcolor[rgb]{0.64,0.00,0.00}{\textbf{#1}}}
\newcommand{\ExtensionTok}[1]{#1}
\newcommand{\FloatTok}[1]{\textcolor[rgb]{0.00,0.00,0.81}{#1}}
\newcommand{\FunctionTok}[1]{\textcolor[rgb]{0.13,0.29,0.53}{\textbf{#1}}}
\newcommand{\ImportTok}[1]{#1}
\newcommand{\InformationTok}[1]{\textcolor[rgb]{0.56,0.35,0.01}{\textbf{\textit{#1}}}}
\newcommand{\KeywordTok}[1]{\textcolor[rgb]{0.13,0.29,0.53}{\textbf{#1}}}
\newcommand{\NormalTok}[1]{#1}
\newcommand{\OperatorTok}[1]{\textcolor[rgb]{0.81,0.36,0.00}{\textbf{#1}}}
\newcommand{\OtherTok}[1]{\textcolor[rgb]{0.56,0.35,0.01}{#1}}
\newcommand{\PreprocessorTok}[1]{\textcolor[rgb]{0.56,0.35,0.01}{\textit{#1}}}
\newcommand{\RegionMarkerTok}[1]{#1}
\newcommand{\SpecialCharTok}[1]{\textcolor[rgb]{0.81,0.36,0.00}{\textbf{#1}}}
\newcommand{\SpecialStringTok}[1]{\textcolor[rgb]{0.31,0.60,0.02}{#1}}
\newcommand{\StringTok}[1]{\textcolor[rgb]{0.31,0.60,0.02}{#1}}
\newcommand{\VariableTok}[1]{\textcolor[rgb]{0.00,0.00,0.00}{#1}}
\newcommand{\VerbatimStringTok}[1]{\textcolor[rgb]{0.31,0.60,0.02}{#1}}
\newcommand{\WarningTok}[1]{\textcolor[rgb]{0.56,0.35,0.01}{\textbf{\textit{#1}}}}
\usepackage{longtable,booktabs,array}
\usepackage{calc} % for calculating minipage widths
% Correct order of tables after \paragraph or \subparagraph
\usepackage{etoolbox}
\makeatletter
\patchcmd\longtable{\par}{\if@noskipsec\mbox{}\fi\par}{}{}
\makeatother
% Allow footnotes in longtable head/foot
\IfFileExists{footnotehyper.sty}{\usepackage{footnotehyper}}{\usepackage{footnote}}
\makesavenoteenv{longtable}
\usepackage{graphicx}
\makeatletter
\newsavebox\pandoc@box
\newcommand*\pandocbounded[1]{% scales image to fit in text height/width
  \sbox\pandoc@box{#1}%
  \Gscale@div\@tempa{\textheight}{\dimexpr\ht\pandoc@box+\dp\pandoc@box\relax}%
  \Gscale@div\@tempb{\linewidth}{\wd\pandoc@box}%
  \ifdim\@tempb\p@<\@tempa\p@\let\@tempa\@tempb\fi% select the smaller of both
  \ifdim\@tempa\p@<\p@\scalebox{\@tempa}{\usebox\pandoc@box}%
  \else\usebox{\pandoc@box}%
  \fi%
}
% Set default figure placement to htbp
\def\fps@figure{htbp}
\makeatother
\setlength{\emergencystretch}{3em} % prevent overfull lines
\providecommand{\tightlist}{%
  \setlength{\itemsep}{0pt}\setlength{\parskip}{0pt}}
\setcounter{secnumdepth}{5}
\usepackage{booktabs}
\usepackage{longtable}
\usepackage{array}
\usepackage{multirow}
\usepackage{wrapfig}
\usepackage{float}
\usepackage{colortbl}
\usepackage{pdflscape}
\usepackage{tabu}
\usepackage{threeparttable}
\usepackage{threeparttablex}
\usepackage[normalem]{ulem}
\usepackage{makecell}
\usepackage{xcolor}
\usepackage{bookmark}
\IfFileExists{xurl.sty}{\usepackage{xurl}}{} % add URL line breaks if available
\urlstyle{same}
\hypersetup{
  pdftitle={Mortalidad por cáncer de pulmón en Ohio (décadas de 1960 a 1980)},
  hidelinks,
  pdfcreator={LaTeX via pandoc}}

\title{Mortalidad por cáncer de pulmón en Ohio (décadas de 1960 a 1980)}
\author{\textbf{Maria Paula Camargo Rincon}\footnote{\href{mailto:mcamargori@unal.edu.co}{\nolinkurl{mcamargori@unal.edu.co}}}
\textbf{Laura Katherine Martínez Castiblanco}\footnote{\href{mailto:laumartinezca@unal.edu.co}{\nolinkurl{laumartinezca@unal.edu.co}}}}
\date{}

\begin{document}
\maketitle

{
\setcounter{tocdepth}{3}
\tableofcontents
}
\section{Introducción}\label{PrInv}

Este conjunto de datos se deriva de un estudio de mortalidad por cáncer
de pulmón en Ohio realizado entre 1969 y 1971 a nivel de condado. La
base de datos \texttt{ohiolung} contiene información para los 88
condados del estado de Ohio, (EE. UU.). La base de datos proviene del
GeoDa (Center for Geospatial Analysis and Computation) y presenta 42
variables. En las estas se presentan datos referentes a la mortalidad
por cáncer de pulmón estratificada por género (masculino/femenino) y
raza (blanca/negra), junto variables geográficas básicas.

La siguiente tabla resume las variables principales de la base de datos:

\begingroup\fontsize{10}{12}\selectfont

\begin{longtable}[t]{>{}l>{\raggedright\arraybackslash}p{10cm}}
\toprule
Variable & Descripcion\\
\midrule
\textbf{\cellcolor{gray!10}{CountyID}} & \cellcolor{gray!10}{ID secuencial del condado (orden alfabético)}\\
\textbf{NAME} & Nombre del condado\\
\textbf{\cellcolor{gray!10}{FIPSNO}} & \cellcolor{gray!10}{Código FIPS del condado (numérico)}\\
\textbf{AREA} & Área del polígono (condado)\\
\textbf{\cellcolor{gray!10}{PERIMETER}} & \cellcolor{gray!10}{Perímetro del polígono (condado)}\\
\addlinespace
\textbf{RECORD\_ID} & ID único del registro\\
\textbf{\cellcolor{gray!10}{COUNTYID}} & \cellcolor{gray!10}{ID del condado}\\
\textbf{LGRyy} & Casos de cáncer de pulmón para género G (M/F) y raza R (W/B) en año yy (1968, 1978, 1988)\\
\textbf{\cellcolor{gray!10}{POPGRyy}} & \cellcolor{gray!10}{Población en riesgo para género G y raza R en año yy}\\
\textbf{LGyy} & Total de casos de cáncer de pulmón por género G en año yy\\
\addlinespace
\textbf{\cellcolor{gray!10}{POPGyy}} & \cellcolor{gray!10}{Población total en riesgo por género G en año yy}\\
\bottomrule
\multicolumn{2}{l}{\rule{0pt}{1em}Tabla 1. Definición de variables principales en la base de datos ohlung.shp.}\\
\end{longtable}
\endgroup{}

Hay 12 variables por año (4 combinaciones de género-raza × 3 tipos:
casos detallados, población detallada, totales por género)

\begingroup\fontsize{9}{11}\selectfont

\begin{longtable}{rlrrrrrr}
\toprule
\textbf{COUNTYID} & \textbf{NAME} & \textbf{FIPSNO} & \textbf{AREA} & \textbf{PERIMETER} & \textbf{LMW68} & \textbf{POPMW68} & \textbf{LMB68}\\
\midrule
\cellcolor{gray!10}{48} & \cellcolor{gray!10}{Lucas} & \cellcolor{gray!10}{39095} & \cellcolor{gray!10}{873382000} & \cellcolor{gray!10}{164533} & \cellcolor{gray!10}{128} & \cellcolor{gray!10}{205421} & \cellcolor{gray!10}{19}\\
26 & Fulton & 39051 & 1054690000 & 134891 & 4 & 15521 & 0\\
\cellcolor{gray!10}{28} & \cellcolor{gray!10}{Geauga} & \cellcolor{gray!10}{39055} & \cellcolor{gray!10}{1005750000} & \cellcolor{gray!10}{146204} & \cellcolor{gray!10}{4} & \cellcolor{gray!10}{30415} & \cellcolor{gray!10}{0}\\
86 & Williams & 39171 & 1089180000 & 136740 & 5 & 16252 & 0\\
\cellcolor{gray!10}{18} & \cellcolor{gray!10}{Cuyahoga} & \cellcolor{gray!10}{39035} & \cellcolor{gray!10}{1242480000} & \cellcolor{gray!10}{173664} & \cellcolor{gray!10}{435} & \cellcolor{gray!10}{677185} & \cellcolor{gray!10}{87}\\
\addlinespace
62 & Ottawa & 39123 & 634883000 & 128356 & 10 & 17638 & 0\\
\bottomrule
\multicolumn{8}{l}{\rule{0pt}{1em}Tabla 2. Primeros 6 datos de la base de datos original.}\\
\end{longtable}
\endgroup{}

\begingroup\fontsize{9}{11}\selectfont

\begin{longtable}{rrrrrrrrrr}
\toprule
\textbf{POPMB68} & \textbf{LM68} & \textbf{POPM68} & \textbf{LFW68} & \textbf{POPFW68} & \textbf{LFB68} & \textbf{POPFB68} & \textbf{LF68} & \textbf{POPF68} & \textbf{LMW78}\\
\midrule
\cellcolor{gray!10}{25956} & \cellcolor{gray!10}{147} & \cellcolor{gray!10}{231377} & \cellcolor{gray!10}{27} & \cellcolor{gray!10}{220106} & \cellcolor{gray!10}{1} & \cellcolor{gray!10}{28527} & \cellcolor{gray!10}{28} & \cellcolor{gray!10}{248633} & \cellcolor{gray!10}{149}\\
50 & 4 & 15571 & 1 & 16358 & 0 & 47 & 1 & 16405 & 10\\
\cellcolor{gray!10}{472} & \cellcolor{gray!10}{4} & \cellcolor{gray!10}{30887} & \cellcolor{gray!10}{3} & \cellcolor{gray!10}{31131} & \cellcolor{gray!10}{0} & \cellcolor{gray!10}{496} & \cellcolor{gray!10}{3} & \cellcolor{gray!10}{31627} & \cellcolor{gray!10}{11}\\
24 & 5 & 16276 & 0 & 17216 & 0 & 20 & 0 & 17236 & 7\\
\cellcolor{gray!10}{157046} & \cellcolor{gray!10}{522} & \cellcolor{gray!10}{834231} & \cellcolor{gray!10}{88} & \cellcolor{gray!10}{734014} & \cellcolor{gray!10}{23} & \cellcolor{gray!10}{176614} & \cellcolor{gray!10}{111} & \cellcolor{gray!10}{910628} & \cellcolor{gray!10}{442}\\
\addlinespace
164 & 10 & 17802 & 0 & 18430 & 0 & 182 & 0 & 18612 & 8\\
\bottomrule
\end{longtable}
\endgroup{}

\begingroup\fontsize{9}{11}\selectfont

\begin{longtable}{rrrrrrrrr}
\toprule
\textbf{POPMW78} & \textbf{LMB78} & \textbf{POPMB78} & \textbf{LM78} & \textbf{POPM78} & \textbf{LFW78} & \textbf{POPFW78} & \textbf{LFB78} & \textbf{POPFB78}\\
\midrule
\cellcolor{gray!10}{194976} & \cellcolor{gray!10}{34} & \cellcolor{gray!10}{30252} & \cellcolor{gray!10}{183} & \cellcolor{gray!10}{225228} & \cellcolor{gray!10}{48} & \cellcolor{gray!10}{210307} & \cellcolor{gray!10}{4} & \cellcolor{gray!10}{33790}\\
18268 & 0 & 75 & 10 & 18343 & 3 & 19080 & 0 & 70\\
\cellcolor{gray!10}{35839} & \cellcolor{gray!10}{1} & \cellcolor{gray!10}{604} & \cellcolor{gray!10}{12} & \cellcolor{gray!10}{36443} & \cellcolor{gray!10}{3} & \cellcolor{gray!10}{36148} & \cellcolor{gray!10}{0} & \cellcolor{gray!10}{619}\\
17678 & 1 & 44 & 8 & 17722 & 0 & 18126 & 0 & 70\\
\cellcolor{gray!10}{567350} & \cellcolor{gray!10}{138} & \cellcolor{gray!10}{161461} & \cellcolor{gray!10}{580} & \cellcolor{gray!10}{728811} & \cellcolor{gray!10}{165} & \cellcolor{gray!10}{623107} & \cellcolor{gray!10}{36} & \cellcolor{gray!10}{186806}\\
\addlinespace
19756 & 0 & 170 & 8 & 19926 & 6 & 20388 & 0 & 188\\
\bottomrule
\end{longtable}
\endgroup{}

\begingroup\fontsize{9}{11}\selectfont

\begin{longtable}{rrrrrrrrr}
\toprule
\textbf{LF78} & \textbf{POPF78} & \textbf{LMW88} & \textbf{POPMW88} & \textbf{LMB88} & \textbf{POPMB88} & \textbf{LM88} & \textbf{POPM88} & \textbf{LFW88}\\
\midrule
\cellcolor{gray!10}{52} & \cellcolor{gray!10}{244097} & \cellcolor{gray!10}{150} & \cellcolor{gray!10}{185367} & \cellcolor{gray!10}{39} & \cellcolor{gray!10}{36104} & \cellcolor{gray!10}{189} & \cellcolor{gray!10}{221471} & \cellcolor{gray!10}{108}\\
3 & 19150 & 8 & 18600 & 1 & 62 & 9 & 18662 & 2\\
\cellcolor{gray!10}{3} & \cellcolor{gray!10}{36767} & \cellcolor{gray!10}{17} & \cellcolor{gray!10}{37759} & \cellcolor{gray!10}{1} & \cellcolor{gray!10}{1051} & \cellcolor{gray!10}{18} & \cellcolor{gray!10}{38810} & \cellcolor{gray!10}{12}\\
0 & 18196 & 10 & 17822 & 0 & 61 & 10 & 17883 & 7\\
\cellcolor{gray!10}{201} & \cellcolor{gray!10}{809913} & \cellcolor{gray!10}{453} & \cellcolor{gray!10}{499655} & \cellcolor{gray!10}{188} & \cellcolor{gray!10}{178432} & \cellcolor{gray!10}{641} & \cellcolor{gray!10}{678087} & \cellcolor{gray!10}{269}\\
\addlinespace
6 & 20576 & 16 & 19164 & 0 & 205 & 16 & 19369 & 9\\
\bottomrule
\end{longtable}
\endgroup{}

\begingroup\fontsize{9}{11}\selectfont

\begin{longtable}{rrrrr}
\toprule
\textbf{POPFW88} & \textbf{LFB88} & \textbf{POPFB88} & \textbf{LF88} & \textbf{POPF88}\\
\midrule
\cellcolor{gray!10}{200458} & \cellcolor{gray!10}{8} & \cellcolor{gray!10}{41115} & \cellcolor{gray!10}{116} & \cellcolor{gray!10}{241573}\\
19444 & 0 & 65 & 2 & 19509\\
\cellcolor{gray!10}{37973} & \cellcolor{gray!10}{0} & \cellcolor{gray!10}{1040} & \cellcolor{gray!10}{12} & \cellcolor{gray!10}{39013}\\
18506 & 0 & 110 & 7 & 18616\\
\cellcolor{gray!10}{552388} & \cellcolor{gray!10}{83} & \cellcolor{gray!10}{207628} & \cellcolor{gray!10}{352} & \cellcolor{gray!10}{760016}\\
\addlinespace
20063 & 1 & 221 & 10 & 20284\\
\bottomrule
\end{longtable}
\endgroup{}

\begin{center}\rule{0.5\linewidth}{0.5pt}\end{center}

\section{Análisis descriptivo}\label{AnalisisDescriptivo}

En esta sección se realiza una exploración de la base de datos con tal
de identificar relaciones entre las variables y establecer la estructura
inicial del modelo de regresión lineal múltiple.

La Tabla 3 se presenta los principales estadísticos descriptivos de los
casos de mortalidad por cáncer de pulmón y las poblaciones en riesgo. Se
observa una alta variabilidad en el número de casos, con desviaciones
estándar frecuentemente superiores a las medias. Los promedios de casos
son consistentemente más altos en hombres que en mujeres y en la
población blanca respecto a la negra. Además, se evidencia un incremento
temporal en el número promedio de casos entre 1968 y 1988.

\begingroup\fontsize{9}{11}\selectfont

\begin{longtable}{lrlr}
\toprule
Variable & Valor & Variable & Valor\\
\midrule
\cellcolor{gray!10}{LMW68 Media} & \cellcolor{gray!10}{26.88} & \cellcolor{gray!10}{LFB78 Min} & \cellcolor{gray!10}{0.00}\\
LMW68 SD & 58.57 & LFB78 Max & 36.00\\
\cellcolor{gray!10}{LMW68 Min} & \cellcolor{gray!10}{1.00} & \cellcolor{gray!10}{LF78 Media} & \cellcolor{gray!10}{13.78}\\
LMW68 Max & 435.00 & LF78 SD & 29.42\\
\cellcolor{gray!10}{LMB68 Media} & \cellcolor{gray!10}{3.26} & \cellcolor{gray!10}{LF78 Min} & \cellcolor{gray!10}{0.00}\\
\addlinespace
LMB68 SD & 11.75 & LF78 Max & 201.00\\
\cellcolor{gray!10}{LMB68 Min} & \cellcolor{gray!10}{0.00} & \cellcolor{gray!10}{LFW88 Media} & \cellcolor{gray!10}{22.33}\\
LMB68 Max & 87.00 & LFW88 SD & 41.47\\
\cellcolor{gray!10}{LM68 Media} & \cellcolor{gray!10}{30.14} & \cellcolor{gray!10}{LFW88 Min} & \cellcolor{gray!10}{0.00}\\
LM68 SD & 70.12 & LFW88 Max & 269.00\\
\addlinespace
\cellcolor{gray!10}{LM68 Min} & \cellcolor{gray!10}{1.00} & \cellcolor{gray!10}{LFB88 Media} & \cellcolor{gray!10}{2.60}\\
LM68 Max & 522.00 & LFB88 SD & 10.49\\
\cellcolor{gray!10}{LMW78 Media} & \cellcolor{gray!10}{35.82} & \cellcolor{gray!10}{LFB88 Min} & \cellcolor{gray!10}{0.00}\\
LMW78 SD & 63.42 & LFB88 Max & 83.00\\
\cellcolor{gray!10}{LMW78 Min} & \cellcolor{gray!10}{3.00} & \cellcolor{gray!10}{LF88 Media} & \cellcolor{gray!10}{24.93}\\
\addlinespace
LMW78 Max & 442.00 & LF88 SD & 51.23\\
\cellcolor{gray!10}{LMB78 Media} & \cellcolor{gray!10}{5.01} & \cellcolor{gray!10}{LF88 Min} & \cellcolor{gray!10}{1.00}\\
LMB78 SD & 17.79 & LF88 Max & 352.00\\
\cellcolor{gray!10}{LMB78 Min} & \cellcolor{gray!10}{0.00} & \cellcolor{gray!10}{POPMW68 Media} & \cellcolor{gray!10}{52983.82}\\
LMB78 Max & 138.00 & POPMW68 SD & 95322.08\\
\addlinespace
\cellcolor{gray!10}{LM78 Media} & \cellcolor{gray!10}{40.83} & \cellcolor{gray!10}{POPMB68 Media} & \cellcolor{gray!10}{5329.95}\\
LM78 SD & 80.55 & POPMB68 SD & 19460.53\\
\cellcolor{gray!10}{LM78 Min} & \cellcolor{gray!10}{3.00} & \cellcolor{gray!10}{POPM68 Media} & \cellcolor{gray!10}{58313.77}\\
LM78 Max & 580.00 & POPM68 SD & 114110.27\\
\cellcolor{gray!10}{LMW88 Media} & \cellcolor{gray!10}{43.01} & \cellcolor{gray!10}{POPMW78 Media} & \cellcolor{gray!10}{53411.91}\\
\addlinespace
LMW88 SD & 68.72 & POPMW78 SD & 85016.85\\
\cellcolor{gray!10}{LMW88 Min} & \cellcolor{gray!10}{2.00} & \cellcolor{gray!10}{POPMB78 Media} & \cellcolor{gray!10}{5915.57}\\
LMW88 Max & 453.00 & POPMB78 SD & 20732.39\\
\cellcolor{gray!10}{LMB88 Media} & \cellcolor{gray!10}{6.22} & \cellcolor{gray!10}{POPM78 Media} & \cellcolor{gray!10}{59327.48}\\
LMB88 SD & 23.11 & POPM78 SD & 104851.91\\
\addlinespace
\cellcolor{gray!10}{LMB88 Min} & \cellcolor{gray!10}{0.00} & \cellcolor{gray!10}{POPMW88 Media} & \cellcolor{gray!10}{52407.83}\\
LMB88 Max & 188.00 & POPMW88 SD & 79338.55\\
\cellcolor{gray!10}{LM88 Media} & \cellcolor{gray!10}{49.23} & \cellcolor{gray!10}{POPMB88 Media} & \cellcolor{gray!10}{6799.70}\\
LM88 SD & 90.40 & POPMB88 SD & 23382.85\\
\cellcolor{gray!10}{LM88 Min} & \cellcolor{gray!10}{2.00} & \cellcolor{gray!10}{POPM88 Media} & \cellcolor{gray!10}{59207.53}\\
\addlinespace
LM88 Max & 641.00 & POPM88 SD & 101568.08\\
\cellcolor{gray!10}{LFW68 Media} & \cellcolor{gray!10}{6.12} & \cellcolor{gray!10}{POPFW68 Media} & \cellcolor{gray!10}{56094.52}\\
LFW68 SD & 12.68 & POPFW68 SD & 102927.37\\
\cellcolor{gray!10}{LFW68 Min} & \cellcolor{gray!10}{0.00} & \cellcolor{gray!10}{POPFB68 Media} & \cellcolor{gray!10}{5862.83}\\
LFW68 Max & 88.00 & POPFB68 SD & 21878.07\\
\addlinespace
\cellcolor{gray!10}{LFB68 Media} & \cellcolor{gray!10}{0.62} & \cellcolor{gray!10}{POPF68 Media} & \cellcolor{gray!10}{61957.35}\\
LFB68 SD & 2.82 & POPF68 SD & 124072.32\\
\cellcolor{gray!10}{LFB68 Min} & \cellcolor{gray!10}{0.00} & \cellcolor{gray!10}{POPFW78 Media} & \cellcolor{gray!10}{56682.57}\\
LFB68 Max & 23.00 & POPFW78 SD & 92507.22\\
\cellcolor{gray!10}{LF68 Media} & \cellcolor{gray!10}{6.75} & \cellcolor{gray!10}{POPFB78 Media} & \cellcolor{gray!10}{6638.98}\\
\addlinespace
LF68 SD & 15.32 & POPFB78 SD & 23892.41\\
\cellcolor{gray!10}{LF68 Min} & \cellcolor{gray!10}{0.00} & \cellcolor{gray!10}{POPF78 Media} & \cellcolor{gray!10}{63321.55}\\
LF68 Max & 111.00 & POPF78 SD & 115408.59\\
\cellcolor{gray!10}{LFW78 Media} & \cellcolor{gray!10}{12.45} & \cellcolor{gray!10}{POPFW88 Media} & \cellcolor{gray!10}{55763.84}\\
LFW78 SD & 24.74 & POPFW88 SD & 86678.27\\
\addlinespace
\cellcolor{gray!10}{LFW78 Min} & \cellcolor{gray!10}{0.00} & \cellcolor{gray!10}{POPFB88 Media} & \cellcolor{gray!10}{7650.48}\\
LFW78 Max & 165.00 & POPFB88 SD & 27083.88\\
\cellcolor{gray!10}{LFB78 Media} & \cellcolor{gray!10}{1.33} & \cellcolor{gray!10}{POPF88 Media} & \cellcolor{gray!10}{63414.32}\\
LFB78 SD & 4.85 & POPF88 SD & 112456.93\\
\bottomrule
\multicolumn{4}{l}{\rule{0pt}{1em}Tabla 3. Estadísticas descriptivas de casos y población.}\\
\end{longtable}
\endgroup{}

Evaluando las distribuciones, se observa una asimetrica a la derecha en
cada grupo de la problación evaluada, dicionalmente se evidencia que los
condados más poblados tienden a presentar menos riesgo de cáncer de
pulmón, muy seguramente por la separacion de zonas urbanas y rurales.

\pandocbounded{\includegraphics[keepaspectratio]{Plantilla_informe2_files/figure-latex/graficos-univariados-1.pdf}}
\pandocbounded{\includegraphics[keepaspectratio]{Plantilla_informe2_files/figure-latex/graficos-univariados-2.pdf}}

\pandocbounded{\includegraphics[keepaspectratio]{Plantilla_informe2_files/figure-latex/graficos-univariados2-1.pdf}}

Las matrices de correlación (presentadas como las combinaciones de :
género y raza en los 3 años, genero femenino en los 3 año y género y
raza únicamente en el último año) muestran una fuerte correlación
positiva entre el número de casos y la población en riesgo
correspondiente, tanto en hombres como en mujeres. Asimismo, se observa
una correlación leve positiva entre los casos y las variables
geográficas \texttt{AREA} y \texttt{PERIMETER}, aunque estas dos últimas
están también correlacionadas entre sí.

A su vez se observa que las variables geográficas aportan información
limitada a la hora de explicar el comportamiento de cualquera de las
variables aya sea de numero de casos o de población en riesgo de cáncer
pulmonar.

\pandocbounded{\includegraphics[keepaspectratio]{Plantilla_informe2_files/figure-latex/correlaciones-1.pdf}}

Por otro lado es de destacar que se observa una alta correlacion lineal
en el tiempo, luego es recomendable añadir un componente que explique
esta temporalidad al modelo, o evaluar la base de datos para tratar el
año como una variable descriptiva para un modelo lineal multiple.

Para cumplir con el requisito de incluir al menos una variable
cualitativa, y teniendo en cuenta las observaciones anteriores, se
reestructuraron los datos al formato largo utilizando exclusivamente las
variables originales de la base de datos (es decir se reorganizó la base
de datos, sin alterar los datos originales). De esta forma se generaron
las variables categóricas \texttt{gender} (Masculino/Femenino) y
\texttt{year} (1968, 1978, 1988).

\begin{longtable}[]{@{}rlrrllrr@{}}
\caption{Vista de los datos en formato largo}\tabularnewline
\toprule\noalign{}
COUNTYID & NAME & AREA & PERIMETER & gender & year & cases & pop \\
\midrule\noalign{}
\endfirsthead
\toprule\noalign{}
COUNTYID & NAME & AREA & PERIMETER & gender & year & cases & pop \\
\midrule\noalign{}
\endhead
\bottomrule\noalign{}
\endlastfoot
48 & Lucas & 873382000 & 164533 & Male & 1968 & 147 & 231377 \\
48 & Lucas & 873382000 & 164533 & Female & 1968 & 28 & 248633 \\
48 & Lucas & 873382000 & 164533 & Male & 1978 & 183 & 225228 \\
48 & Lucas & 873382000 & 164533 & Female & 1978 & 52 & 244097 \\
48 & Lucas & 873382000 & 164533 & Male & 1988 & 189 & 221471 \\
48 & Lucas & 873382000 & 164533 & Female & 1988 & 116 & 241573 \\
26 & Fulton & 1054690000 & 134891 & Male & 1968 & 4 & 15571 \\
26 & Fulton & 1054690000 & 134891 & Female & 1968 & 1 & 16405 \\
26 & Fulton & 1054690000 & 134891 & Male & 1978 & 10 & 18343 \\
26 & Fulton & 1054690000 & 134891 & Female & 1978 & 3 & 19150 \\
\end{longtable}

Considerando la nueva estructura de los datos se realiza un análisis
adicional sobre la variable respuesta \texttt{cases} para evaluar su
comportamiento en función del género y el año.

\pandocbounded{\includegraphics[keepaspectratio]{Plantilla_informe2_files/figure-latex/boxplots-1.pdf}}
\pandocbounded{\includegraphics[keepaspectratio]{Plantilla_informe2_files/figure-latex/boxplots-2.pdf}}

Se observa una precencia de datos atípicos por años, mas dispersión
conforme pasan los años, sin embargo la media parece mantenerse.Más
dispersos en hombres que mujeres pero media constante, más datos
atípicos por hombres que por mujeres.

\pandocbounded{\includegraphics[keepaspectratio]{Plantilla_informe2_files/figure-latex/scatterplot-1.pdf}}

\begin{center}\rule{0.5\linewidth}{0.5pt}\end{center}

\section{Modelo y verificación de supuestos}\label{ejem}

Considerando el aálisis anterior se plantea el siguiente modelo:

\[ \text{cases} = \beta_0 + \beta_1 \text{AREA} + \beta_2 \text{PERIMETER} + \beta_3 \text{gender} + \beta_4 \text{year} + \beta_5 \text{pop} + \epsilon \]

donde \texttt{cases} corresponde a las variables \texttt{LMyy} o
\texttt{LFyy} según el género y año, y \texttt{pop} es la población en
riesgo correspondiente.

Los resultados de este modelo muestran que la población (\texttt{pop})
es altamente significativa y tiene el mayor impacto, como era esperado.
El género masculino se asocia con un promedio de aproximadamente 26.8
casos adicionales respecto al femenino (controlando por las demás
variables). Además, se observa un aumento significativo en el número de
casos en 1978 y 1988 respecto a 1968.

\begin{verbatim}
## 
## Call:
## lm(formula = cases ~ AREA + PERIMETER + gender + year + pop, 
##     data = ohlung_long)
## 
## Residuals:
##      Min       1Q   Median       3Q      Max 
## -302.741   -9.193   -0.937    9.265  293.559 
## 
## Coefficients:
##               Estimate Std. Error t value Pr(>|t|)    
## (Intercept) -3.308e+01  1.209e+01  -2.736  0.00643 ** 
## AREA        -9.100e-09  7.592e-09  -1.199  0.23124    
## PERIMETER    1.312e-04  9.966e-05   1.316  0.18862    
## genderMale   2.680e+01  2.726e+00   9.831  < 2e-16 ***
## year1978     8.295e+00  3.338e+00   2.485  0.01326 *  
## year1988     1.807e+01  3.338e+00   5.415 9.36e-08 ***
## pop          4.781e-04  1.228e-05  38.928  < 2e-16 ***
## ---
## Signif. codes:  0 '***' 0.001 '**' 0.01 '*' 0.05 '.' 0.1 ' ' 1
## 
## Residual standard error: 31.31 on 521 degrees of freedom
## Multiple R-squared:  0.7613, Adjusted R-squared:  0.7586 
## F-statistic:   277 on 6 and 521 DF,  p-value: < 2.2e-16
\end{verbatim}

Las variables \texttt{AREA} y \texttt{PERIMETER} no resultan
estadísticamente significativas. Un modelo reducido sin estas variables
geográficas presenta un \(R^2\) prácticamente idéntico (0.7605 vs
0.7613), lo que confirma su escaso aporte y se propone en cambio el
siguiente modelo reducido:

\[ \text{cases} = \beta_0  + \beta_1 \text{gender} + \beta_2 \text{year} + \beta_3 \text{pop} + \epsilon \]

\begin{verbatim}
## 
## Call:
## lm(formula = cases ~ gender + year + pop, data = ohlung_long)
## 
## Residuals:
##      Min       1Q   Median       3Q      Max 
## -301.984   -9.110   -1.014    9.029  294.680 
## 
## Coefficients:
##               Estimate Std. Error t value Pr(>|t|)    
## (Intercept) -2.380e+01  2.828e+00  -8.415 3.77e-16 ***
## genderMale   2.680e+01  2.726e+00   9.834  < 2e-16 ***
## year1978     8.293e+00  3.338e+00   2.485   0.0133 *  
## year1988     1.807e+01  3.338e+00   5.415 9.36e-08 ***
## pop          4.796e-04  1.220e-05  39.314  < 2e-16 ***
## ---
## Signif. codes:  0 '***' 0.001 '**' 0.01 '*' 0.05 '.' 0.1 ' ' 1
## 
## Residual standard error: 31.31 on 523 degrees of freedom
## Multiple R-squared:  0.7605, Adjusted R-squared:  0.7586 
## F-statistic: 415.1 on 4 and 523 DF,  p-value: < 2.2e-16
\end{verbatim}

Además de esto, hemos realizado regresión por ridge y lasso mostradas a
continuación:

\begin{Shaded}
\begin{Highlighting}[]
\NormalTok{X }\OtherTok{\textless{}{-}} \FunctionTok{as.matrix}\NormalTok{(ohlung\_long }\SpecialCharTok{\%\textgreater{}\%}\NormalTok{ dplyr}\SpecialCharTok{::}\FunctionTok{select}\NormalTok{(}\SpecialCharTok{{-}}\NormalTok{cases))}
\NormalTok{y }\OtherTok{\textless{}{-}}\NormalTok{ ohlung\_long}\SpecialCharTok{$}\NormalTok{cases}
\NormalTok{ridge\_model }\OtherTok{\textless{}{-}} \FunctionTok{glmnet}\NormalTok{(X, y, }\AttributeTok{alpha =} \DecValTok{0}\NormalTok{)}
\FunctionTok{print}\NormalTok{(ridge\_model)}
\end{Highlighting}
\end{Shaded}

\begin{verbatim}
## 
## Call:  glmnet(x = X, y = y, alpha = 0) 
## 
##     Df  %Dev Lambda
## 1    5  0.00  53370
## 2    5  0.19  48630
## 3    5  0.21  44310
## 4    5  0.23  40370
## 5    5  0.25  36790
## 6    5  0.28  33520
## 7    5  0.31  30540
## 8    5  0.33  27830
## 9    5  0.37  25350
## 10   5  0.40  23100
## 11   5  0.44  21050
## 12   5  0.48  19180
## 13   5  0.53  17480
## 14   5  0.58  15920
## 15   5  0.64  14510
## 16   5  0.70  13220
## 17   5  0.77  12050
## 18   5  0.84  10980
## 19   5  0.93  10000
## 20   5  1.01   9112
## 21   5  1.11   8303
## 22   5  1.22   7565
## 23   5  1.34   6893
## 24   5  1.46   6281
## 25   5  1.61   5723
## 26   5  1.76   5214
## 27   5  1.93   4751
## 28   5  2.11   4329
## 29   5  2.31   3944
## 30   5  2.53   3594
## 31   5  2.77   3275
## 32   5  3.03   2984
## 33   5  3.32   2719
## 34   5  3.63   2477
## 35   5  3.96   2257
## 36   5  4.33   2057
## 37   5  4.73   1874
## 38   5  5.17   1707
## 39   5  5.64   1556
## 40   5  6.16   1418
## 41   5  6.71   1292
## 42   5  7.32   1177
## 43   5  7.97   1072
## 44   5  8.67    977
## 45   5  9.43    890
## 46   5 10.25    811
## 47   5 11.13    739
## 48   5 12.07    673
## 49   5 13.08    614
## 50   5 14.16    559
## 51   5 15.30    509
## 52   5 16.53    464
## 53   5 17.82    423
## 54   5 19.19    385
## 55   5 20.63    351
## 56   5 22.14    320
## 57   5 23.73    292
## 58   5 25.38    266
## 59   5 27.10    242
## 60   5 28.87    220
## 61   5 30.70    201
## 62   5 32.58    183
## 63   5 34.49    167
## 64   5 36.43    152
## 65   5 38.39    138
## 66   5 40.35    126
## 67   5 42.32    115
## 68   5 44.27    105
## 69   5 46.19     95
## 70   5 48.08     87
## 71   5 49.92     79
## 72   5 51.70     72
## 73   5 53.41     66
## 74   5 55.05     60
## 75   5 56.61     55
## 76   5 58.08     50
## 77   5 59.46     45
## 78   5 60.74     41
## 79   5 61.93     38
## 80   5 63.03     34
## 81   5 64.03     31
## 82   5 64.95     28
## 83   5 65.78     26
## 84   5 66.52     24
## 85   5 67.19     22
## 86   5 67.78     20
## 87   5 68.31     18
## 88   5 68.78     16
## 89   5 69.18     15
## 90   5 69.54     14
## 91   5 69.86     12
## 92   5 70.13     11
## 93   5 70.36     10
## 94   5 70.57      9
## 95   5 70.74      8
## 96   5 70.89      8
## 97   5 71.02      7
## 98   5 71.13      6
## 99   5 71.22      6
## 100  5 71.30      5
\end{verbatim}

\begin{Shaded}
\begin{Highlighting}[]
\NormalTok{cv\_ridge }\OtherTok{\textless{}{-}} \FunctionTok{cv.glmnet}\NormalTok{(X, y, }\AttributeTok{alpha =} \DecValTok{0}\NormalTok{)}

\NormalTok{best\_lambda }\OtherTok{\textless{}{-}}\NormalTok{ cv\_ridge}\SpecialCharTok{$}\NormalTok{lambda.min}
\NormalTok{best\_lambda}
\end{Highlighting}
\end{Shaded}

\begin{verbatim}
## [1] 5.336935
\end{verbatim}

\begin{Shaded}
\begin{Highlighting}[]
\NormalTok{ridge\_best }\OtherTok{\textless{}{-}} \FunctionTok{glmnet}\NormalTok{(X, y, }\AttributeTok{alpha =} \DecValTok{0}\NormalTok{, }\AttributeTok{lambda =}\NormalTok{ best\_lambda)}

\FunctionTok{coef}\NormalTok{(ridge\_best)}
\end{Highlighting}
\end{Shaded}

\begin{verbatim}
## 8 x 1 sparse Matrix of class "dgCMatrix"
##                        s0
## (Intercept) -1.662474e+03
## COUNTYID    -4.783533e-02
## NAME         .           
## AREA        -7.848866e-09
## PERIMETER    1.378861e-04
## gender       .           
## year         8.359870e-01
## pop          4.382109e-04
\end{verbatim}

La regresión ridge mostro que las variables que tienen un impacto casi
nulo en la variable respuesta son: `gender' y `COUNTYID'.

\begin{Shaded}
\begin{Highlighting}[]
\NormalTok{lasso\_model }\OtherTok{\textless{}{-}} \FunctionTok{glmnet}\NormalTok{(X, y, }\AttributeTok{alpha =} \DecValTok{1}\NormalTok{)}
\end{Highlighting}
\end{Shaded}

\begin{verbatim}
## Warning in storage.mode(x) <- "double": NAs introducidos por coerción
\end{verbatim}

\begin{Shaded}
\begin{Highlighting}[]
\FunctionTok{print}\NormalTok{(lasso\_model)}
\end{Highlighting}
\end{Shaded}

\begin{verbatim}
## 
## Call:  glmnet(x = X, y = y, alpha = 1) 
## 
##    Df  %Dev Lambda
## 1   0  0.00 53.370
## 2   1 11.93 48.630
## 3   1 21.84 44.310
## 4   1 30.06 40.370
## 5   1 36.89 36.790
## 6   1 42.55 33.520
## 7   1 47.26 30.540
## 8   1 51.17 27.830
## 9   1 54.41 25.350
## 10  1 57.10 23.100
## 11  1 59.34 21.050
## 12  1 61.19 19.180
## 13  1 62.74 17.480
## 14  1 64.01 15.920
## 15  1 65.08 14.510
## 16  1 65.96 13.220
## 17  1 66.69 12.050
## 18  1 67.30 10.980
## 19  1 67.80 10.000
## 20  1 68.22  9.112
## 21  1 68.57  8.303
## 22  1 68.86  7.565
## 23  2 69.28  6.893
## 24  2 69.68  6.281
## 25  2 70.00  5.723
## 26  2 70.28  5.214
## 27  2 70.50  4.751
## 28  2 70.69  4.329
## 29  2 70.85  3.944
## 30  2 70.98  3.594
## 31  2 71.09  3.275
## 32  2 71.18  2.984
## 33  2 71.25  2.719
## 34  2 71.31  2.477
## 35  2 71.36  2.257
## 36  2 71.41  2.057
## 37  2 71.44  1.874
## 38  2 71.47  1.707
## 39  2 71.49  1.556
## 40  2 71.51  1.418
## 41  2 71.53  1.292
## 42  2 71.55  1.177
## 43  3 71.56  1.072
## 44  4 71.58  0.977
## 45  4 71.59  0.890
## 46  4 71.60  0.811
## 47  5 71.62  0.739
## 48  5 71.64  0.673
## 49  5 71.65  0.614
## 50  5 71.66  0.559
## 51  5 71.68  0.509
## 52  5 71.68  0.464
## 53  5 71.69  0.423
## 54  5 71.70  0.385
## 55  5 71.70  0.351
## 56  5 71.71  0.320
## 57  5 71.71  0.292
## 58  5 71.71  0.266
## 59  5 71.72  0.242
## 60  5 71.72  0.220
## 61  5 71.72  0.201
## 62  5 71.72  0.183
## 63  5 71.72  0.167
## 64  5 71.72  0.152
## 65  5 71.72  0.138
## 66  5 71.72  0.126
\end{verbatim}

\begin{Shaded}
\begin{Highlighting}[]
\NormalTok{cv\_lasso }\OtherTok{\textless{}{-}} \FunctionTok{cv.glmnet}\NormalTok{(X, y, }\AttributeTok{alpha =} \DecValTok{1}\NormalTok{)}

\NormalTok{(best\_lambda\_lasso }\OtherTok{\textless{}{-}}\NormalTok{ cv\_ridge}\SpecialCharTok{$}\NormalTok{lambda.min)}
\end{Highlighting}
\end{Shaded}

\begin{verbatim}
## [1] 5.336935
\end{verbatim}

\begin{Shaded}
\begin{Highlighting}[]
\NormalTok{lasso\_best }\OtherTok{\textless{}{-}} \FunctionTok{glmnet}\NormalTok{(X, y, }\AttributeTok{alpha =} \DecValTok{1}\NormalTok{, }\AttributeTok{lambda =}\NormalTok{ best\_lambda\_lasso)}

\FunctionTok{coef}\NormalTok{(lasso\_best)}
\end{Highlighting}
\end{Shaded}

\begin{verbatim}
## 8 x 1 sparse Matrix of class "dgCMatrix"
##                        s0
## (Intercept) -4.988451e+02
## COUNTYID     .           
## NAME         .           
## AREA         .           
## PERIMETER    .           
## gender       .           
## year         2.529125e-01
## pop          4.299486e-04
\end{verbatim}

Por otro lado el modelo de regresión lasso indica que las variables que
son influyentes son: `year' y `pop'. Por lo que en el modelo final se ha
tomado la decisión de tener como variables explicativas: `gender',`year'
y `pop'.

\[y_{i}=\beta_0+\beta_1x_{i1}+\beta_2x_{i2}+\beta_3x_{i3}+\beta_4x_{i4}\]
donde: - \(x_{i1}\) es la variable dummy referente al género masculino,
es decir, 1 si es masculino 0 en otro caso. - \(x_{i2}\) es la variable
dummy referente al año de 1978. - \(x_{i3}\) es la variable dummy
referente al año de 1988. - \(x_{i4}\) es la variable númerica de la
población en riesgo.

\section{Verificación de supuestos}\label{VerifSupuestos}

Para un primer avistamiento sobre la validación del modelo veamos el
comportamiento de los residuales:

\begin{Shaded}
\begin{Highlighting}[]
\FunctionTok{plot}\NormalTok{(}\FunctionTok{fitted.values}\NormalTok{(modelo\_2),}\FunctionTok{residuals.lm}\NormalTok{(modelo\_2))}
\end{Highlighting}
\end{Shaded}

\pandocbounded{\includegraphics[keepaspectratio]{Plantilla_informe2_files/figure-latex/supuesto1-1.pdf}}
Podemos ver que hay un problema evidente de heteroscedastiidad que más
adelante será confirmado con las respectivas pruebas.

\subsection{\texorpdfstring{\emph{Independencia:}}{Independencia:}}\label{independencia}

\begin{Shaded}
\begin{Highlighting}[]
\FunctionTok{acf}\NormalTok{(}\FunctionTok{residuals.lm}\NormalTok{(modelo\_2))}
\end{Highlighting}
\end{Shaded}

\pandocbounded{\includegraphics[keepaspectratio]{Plantilla_informe2_files/figure-latex/supuesto2-1.pdf}}
Parecen haber algunas correlaciones entre los residuales.

\begin{Shaded}
\begin{Highlighting}[]
\FunctionTok{dwtest}\NormalTok{(modelo\_2)}
\end{Highlighting}
\end{Shaded}

\begin{verbatim}
## 
##  Durbin-Watson test
## 
## data:  modelo_2
## DW = 3.3245, p-value = 1
## alternative hypothesis: true autocorrelation is greater than 0
\end{verbatim}

Debido a que su p valor fue mayor al nivel de significancia (0.05), nos
indica que no se rechaza \(H_0\) no hay autocorrelación en los
residuales, aunque el ACF sugiere algunas correlaciones en ciertos
rezagos, la prueba de Durbin--Watson no encuentra evidencia estadística
suficiente para rechazar la hipótesis de independencia al 5\%. \#\#
\emph{Normalidad:}

\begin{Shaded}
\begin{Highlighting}[]
\FunctionTok{hist}\NormalTok{(}\FunctionTok{residuals.lm}\NormalTok{(modelo\_2),}\AttributeTok{breaks =} \DecValTok{50}\NormalTok{)}
\end{Highlighting}
\end{Shaded}

\pandocbounded{\includegraphics[keepaspectratio]{Plantilla_informe2_files/figure-latex/supuesto4-1.pdf}}

\begin{Shaded}
\begin{Highlighting}[]
\FunctionTok{qqnorm}\NormalTok{(}\FunctionTok{residuals.lm}\NormalTok{(modelo\_2))}
\FunctionTok{qqline}\NormalTok{(}\FunctionTok{residuals.lm}\NormalTok{(modelo\_2))}
\end{Highlighting}
\end{Shaded}

\pandocbounded{\includegraphics[keepaspectratio]{Plantilla_informe2_files/figure-latex/supuesto4-2.pdf}}
Con este histograma no podemos asegurar que la distribución de los
residuales sea una normal, sin embargo se ve que es una distribución
simétrica, y con media en 0. Por otro lado el qqplot nos muestra que
puede presentar problemas de normalidad.

\begin{Shaded}
\begin{Highlighting}[]
\FunctionTok{shapiro.test}\NormalTok{(}\FunctionTok{residuals.lm}\NormalTok{(modelo\_2))}
\end{Highlighting}
\end{Shaded}

\begin{verbatim}
## 
##  Shapiro-Wilk normality test
## 
## data:  residuals.lm(modelo_2)
## W = 0.59974, p-value < 2.2e-16
\end{verbatim}

Con un p valor menor al nivel de significancia (0.05) se rechaza la
\(H_0\) la población viene de una distribución normal.

\subsection{\texorpdfstring{\emph{Hetercedasticidad:}}{Hetercedasticidad:}}\label{hetercedasticidad}

Para evaluar la heteroscedasticidad formalmente, tomando en cuenta que
la población no sigue o proviene de una distrincuón normal, se decidió
aplicar el test de Breusch-pagan.

\begin{Shaded}
\begin{Highlighting}[]
\FunctionTok{bptest}\NormalTok{(modelo\_2)}
\end{Highlighting}
\end{Shaded}

\begin{verbatim}
## 
##  studentized Breusch-Pagan test
## 
## data:  modelo_2
## BP = 240.36, df = 4, p-value < 2.2e-16
\end{verbatim}

\subsection{\texorpdfstring{Con un p valor menor al nivel de
significancia (0.05) se rechaza la \(H_0\) igualdad de varianzas, es
algo que ya se podría detectar a simple vista en el gráfico de los
residuales y los valores
predichos.}{Con un p valor menor al nivel de significancia (0.05) se rechaza la H\_0 igualdad de varianzas, es algo que ya se podría detectar a simple vista en el gráfico de los residuales y los valores predichos.}}\label{con-un-p-valor-menor-al-nivel-de-significancia-0.05-se-rechaza-la-h_0-igualdad-de-varianzas-es-algo-que-ya-se-podruxeda-detectar-a-simple-vista-en-el-gruxe1fico-de-los-residuales-y-los-valores-predichos.}

\section{Concluciones}\label{concluciones}

\begin{itemize}
\tightlist
\item
  La regre
\end{itemize}

\section{Revisión bibliográfica}\label{revisiuxf3n-bibliogruxe1fica}

\textbf{Fuentes propuestas por los consultantes:}\\
- @gareth2021introduction\\
- @pena2013analisis\\
- @schoenberg1935remarks\\
- @corradino1990proximity

\textbf{Fuentes consultadas por el grupo:}\\
- @bib2\\
- @lapointe1994classification

\begin{center}\rule{0.5\linewidth}{0.5pt}\end{center}

\section{Bibliografía}\label{bibliografuxeda}

\end{document}
